{\color{ink}\bfseries 例 4-12 \(\cdot \) 第3问 \(\cdot \) 合订本170 \(\cdot \) 图2}\par \textbf{210.2.4.0/\allowbreak{}24，直连R16}\par\vspace{2bp}{\setlength{\tabcolsep}{1.2bp}\renewcommand{\arraystretch}{1.18}\begin{tabular}{@{}>{\raggedright\arraybackslash}p{39.243bp}>{\raggedright\arraybackslash}p{21.802bp}>{\raggedright\arraybackslash}p{21.802bp}>{\raggedright\arraybackslash}p{21.802bp}>{\raggedright\arraybackslash}p{21.802bp}@{}}
\toprule
\textbf{路由器} & \textbf{初始} & \textbf{0 s} & \textbf{30 s} & \textbf{60 s}\\\midrule
R11 & \(\infty \) & \(\infty \) & \(\infty \) & 4\\
R12 & \(\infty \) & \(\infty \) & \(\infty \) & 4\\
R13 & \(\infty \) & \(\infty \) & 3 & 3\\
R14 & \(\infty \) & \(\infty \) & 3 & 3\\
R15 & \(\infty \) & 2 & 2 & 2\\
R16 & 1 & 1 & 1 & 1\\
\bottomrule\end{tabular}}\par \(\infty \)表示RIP不可达（16）；直连网距离1。首次全体获知：60 s。
